\section{Key Technical Decisions}
\label{sec:tech-decisions}

Section~\ref{sec:timeline} presented the project chronologically. This
section regroups the material by theme, for a reader trying to understand
one subsystem (e.g.\ ``how does grasp geometry work'') rather than one
session.

\subsection{Perception and calibration}

\paragraph{ChArUco extrinsic calibration.} The camera-to-base transform is
computed as $T_{\text{base,cam}} = T_{\text{base,board}} \cdot
T_{\text{cam,board}}^{-1}$, where $T_{\text{cam,board}}$ is estimated live
each frame via \texttt{solvePnP} and $T_{\text{base,board}}$ is a fixed,
manually-measured pose (\texttt{board\_in\_base\_\{x,y,z,roll,pitch,yaw\}}
in \texttt{config.py}). When the board is not visible, the deprojection
service emits a warning and falls back to returning coordinates in the
camera frame rather than failing silently -- a deliberate choice so that a
temporarily-occluded board produces an obviously-wrong (camera-frame)
answer instead of a plausible-looking wrong (stale-base-frame) one.

\paragraph{Handedness (\texttt{charuco\_z\_sign}).} The board pose estimate
can resolve to either a right-handed or left-handed triad depending on
corner-ordering ambiguity; \texttt{charuco\_z\_sign=1.0} (right-handed, Z
toward the camera) is correct and is now the default -- it defaulted to
$-1.0$ early in the project and was corrected once the ambiguity was
understood.

\paragraph{\texttt{z\_offset\_m}.} A small (default 1\,cm) offset is added to
every deprojected point before it reaches the planner, lifting object
coordinates slightly above the measured table surface. This exists because
depth noise and calibration error otherwise occasionally place a computed
pick point fractionally \emph{below} the physical table, which the executor
cannot safely reach.

\paragraph{The VLM pixel-coordinate convention bug (Session 4).} The single
highest-leverage perception fix in the project: every VLM prompt asked for
pixel coordinates in \texttt{[h, w]} (row, column) order, but Qwen-family
models ignore that instruction and always emit their native \texttt{[x, y]}
(column, row) convention. This produced a systematic axis swap between where
the model thought it was pointing and where the deprojection actually
sampled -- a bug that looks like random grounding noise until the swap
pattern is noticed. It was standardized to \texttt{[x, y]} everywhere
(prompts, schemas, deprojection code) rather than fought at the prompt level,
since the model's behavior was consistent, just not the one requested.

\subsection{Grasp and release geometry}

This is the most iterated-on subsystem in the project, because it is where
small, silent unit-convention mismatches directly become millimeter-scale
hardware errors. See \texttt{docs/GRASP\_GEOMETRY\_PIPELINE.md} for the full
derivation; the key decisions:

\begin{itemize}[nosep]
  \item \textbf{Mid-body pick descent.} The gripper descends to
    $\text{descent} = \min(0.5 \times \text{height}, \text{max\_descent})$,
    not to the object's top surface -- this centers the grasp contact point
    on typical rigid objects rather than pinching right at the top edge.
  \item \textbf{Width-aware TCP offset.} The RG2's fingers pivot, so
    flange-to-contact distance is a function of how far apart the fingers
    are, which depends on object width. A piecewise-linear calibration table
    (\texttt{grasp\_geometry.py}) interpolates the correct TCP Z offset per
    pick, reset to a default after release.
  \item \textbf{Release-lift height must match the actual grasp point, not
    the object's full height.} The recurring bug pattern across two
    sessions: release-lift amount was set to the object's \emph{full}
    height, implicitly assuming the object was grasped at its very top --
    but the paired pick descended mid-body. Net effect: every release
    overshot upward by roughly half the object's height. Fixed (Session 4)
    by carrying \texttt{height - descent} (true grasp-point-to-bottom
    distance) as the release-lift value instead of the full height.
  \item \textbf{Double-counted stacking height (VLM\_LLM hybrid, Session 4).}
    A depth-measured true top surface was having the VLM's separate, blind
    size estimate added on top of it a second time by the LLM's own
    stacking formula (\texttt{position.z + size[2]}). Fixed by deriving
    target height from depth and resetting \texttt{position.z} to a base
    reference so the existing formula reconstructs the true top surface
    instead of double-counting.
  \item \textbf{Deterministic LLM-mode geometry (Session 5).} The same class
    of bug recurred in LLM mode, where every prompt asks the model to
    compute \texttt{object\_height} and \texttt{release\_position.z} via its
    own arithmetic -- consistently wrong on real hardware. Rather than
    iterating the prompt again, this was pulled out of the model entirely:
    \texttt{\_fix\_object\_height} recomputes it from the paired pick's
    actual grasp Z minus table surface Z; \texttt{\_fix\_release\_height}
    matches a release against the scene by \emph{footprint containment}
    (bounding box, not point-distance -- a tray is tens of centimeters wide,
    so a release can land validly inside it while being far from its
    registered center point) and nudges only positions that drifted into a
    zone without exactly echoing a registered target/stacking position
    (to avoid re-nudging a deliberate stack or placement).
  \item \textbf{Grasp-width sampling (Session 5).} Deriving \texttt{grasp\_width}
    by deprojecting a bounding box's left/right edge pixels sampled directly
    on the object's silhouette boundary -- the single worst place for valid
    depth (occlusion/parallax) -- and one failed edge sample aborted the
    entire batched deprojection call, crashing an otherwise-valid plan.
    Fixed by sampling inset from the edge and isolating the failure so it
    degrades to ``keep the VLM's own width guess'' instead of crashing.
  \item \textbf{Bounding-box target schema (ROBOAI-19, September).}
    \texttt{targets.*} in \texttt{scene\_mock.json} replaced
    \texttt{position} (center) + \texttt{size} (dimensions) with an explicit
    \texttt{bounds: \{x, y, z\}} of \texttt{[min, max]} pairs, where
    \texttt{bounds.z[1]} is the zone's top surface. Removes an assumed,
    undocumented convention (\texttt{position} = center, half-extent =
    \texttt{size}/2) this project's bugs had repeatedly tripped on, and
    matches how a physical zone is actually measured by hand -- two
    corners, not a center and a half-width. \texttt{objects.*} (things to
    pick) deliberately kept \texttt{position}+\texttt{size}: a pick needs a
    single contact point, not an area. Two hardware-found bugs surfaced
    while wiring this through the LLM-mode prompts: four of six reasoning
    methods (\texttt{cot\_sc}, \texttt{tot}, \texttt{always\_act},
    \texttt{self\_refine}) never had a target-geometry instruction at all,
    which was invisible under the old schema (the model could just copy a
    ready-made \texttt{position}) and became a real hazard under
    \texttt{bounds} (the model must compute a midpoint, and without being
    told to, doesn't reliably land inside the zone).
  \item \textbf{Zone-bounds-aware release spacing (ROBOAI-18, September).}
    \texttt{distribute\_zone\_releases} (\texttt{schemas.py}) used to space
    out colliding releases along an unbounded line/grid, with no notion of
    how big the target zone actually was -- nothing stopped a fourth or
    fifth object from landing just outside the tray's physical edge, and
    \texttt{PlanValidator} only checks global workspace limits, not
    per-zone bounds, so the overflow was caught nowhere. Now takes an
    optional \texttt{targets} registry (the ROBOAI-19 \texttt{bounds}
    schema) and, when a release cluster's anchor matches a known zone,
    lays out occupants on a grid centered on the anchor and clamped to that
    zone's bounds; falls back to the old unbounded behavior when no zone
    matches (the case for point-grounded VLM releases, which have no
    footprint to match against).
  \item \textbf{Multi-object sequential-placement collision (identified,
    not yet fixed).} The September benchmark's dominant failure mode
    (Section~\ref{sec:benchmark}): spacing releases inside a zone's bounds
    keeps \emph{final} positions from overlapping, but nothing prevents the
    arm from clipping an already-placed object while approaching or
    releasing the next one, or a released object rolling into a neighbor
    from a few cm of height error. Twenty-two of twenty-two non-empty
    operator notes in the September rerun describe exactly this, in both
    the LLM and VLM arms -- a real, open problem distinct from the
    zone-bounds fix above, not a claimed contribution.
\end{itemize}

\subsection{Reasoning-method robustness}

Across six independently-implemented reasoning methods (FHP/FFHP, ReAct,
CoT-SC, Self-Refine, Tree-of-Thought, Always-Act), the same categories of
failure recurred and were each fixed once at the shared base-class or helper
level rather than six times:

\begin{itemize}[nosep]
  \item \textbf{Markdown-fenced JSON} (\texttt{```json ... ```}) breaking
    \texttt{json.loads()} -- fixed with a shared
    \texttt{\_strip\_json\_fence()} classmethod.
  \item \textbf{Unclosed \texttt{<think>} blocks} on token-budget exhaustion
    silently producing a non-blank-looking response that skipped the
    existing blank-response retry path -- fixed by extending the
    think-block-stripping regex to also match an \emph{un}closed block.
  \item \textbf{Non-dict actions in a parsed plan} causing an unhandled
    \texttt{TypeError} deep in Pydantic validation instead of a catchable
    planning error -- fixed with an \texttt{isinstance} guard raising a
    dedicated \texttt{ActionParsingError}.
  \item \textbf{Sequential CoT-SC sampling} multiplying a single slow-call
    risk by $k$ -- fixed by parallelizing the $k$ independent samples with a
    thread pool (and adding a lock around now-concurrent metrics updates).
  \item \textbf{No provider request timeout anywhere} -- every provider
    client gained a configurable, bounded request timeout
    (\texttt{Settings.REQUEST\_TIMEOUT\_S}), closing off a class of ``the plan
    was generated but the app already gave up waiting'' failures.
\end{itemize}

The pattern across all of these: \textbf{push correctness fixes down into a
shared base class or helper} rather than into each reasoning method's
individual file, so future reasoning methods inherit the fix automatically.

\subsection{Concurrency and process supervision}

\begin{itemize}[nosep]
  \item \textbf{\texttt{ReentrantCallbackGroup} + \texttt{MultiThreadedExecutor}
    everywhere a node makes nested service calls} (Section~\ref{sec:architecture}) --
    the single structural rule that prevents deadlock in this codebase.
  \item \textbf{The UR driver is treated as inherently flaky, not fixed.}
    Rather than debugging the stock \texttt{ur\_control.launch.py} driver's
    intermittent startup failure, \texttt{UrDriverSupervisor} wraps it with
    retry-with-backoff against an explicit readiness probe, and later
    (Session 4) fast-fails on an \texttt{[ERROR]} log marker instead of
    always waiting out a fixed timeout.
  \item \textbf{Hard guard against duplicate executors.} A leftover
    \texttt{/execute\_skill} action server from a crashed or improperly
    stopped previous session would silently double-run every skill if a new
    stack started on top of it. \texttt{StackSupervisor} checks for this
    explicitly (not just by process name) before allowing a new start.
\end{itemize}

\subsection{Configuration consolidation}

Settings accumulated in three places over the project's history: scattered
per-class Python constants, launch-file default arguments, and
\texttt{config.py}. Two consolidation passes (Session 2's ChArUco defaults
sync, Session 4's full settings consolidation) moved everything into
\texttt{config.py} as the single \texttt{pydantic-settings} source of truth,
each field overridable via a \texttt{ROBOREASON\_}-prefixed environment
variable. The recurring failure mode this fixes is a launch file's hardcoded
default silently drifting out of sync with the authoritative value elsewhere
(e.g.\ Session 2's board-geometry mismatch, Session 3's stale model-name
defaults after a registry refresh) -- consolidation does not prevent drift by
itself, but it reduces the number of places drift can occur from three to
one.
